\section{Related Work}
\label{sec:related}

\paragraph{Performance-measurement and benchmarking methodology}
Performance-evaluation practice emphasizes construct validity, measurement semantics, and protocol
transparency: reported rankings depend on what is counted, at what layer, and under what workload
definition \citep{jain1991,lilja2000,georges2007}. In budgeted LLM inference, the analogous
construct-validity issue is whether a nominal action budget measures the resource whose performance
claim is being made. We therefore separate the nominal budget from successful completions, HTTP
retries, generated tokens, observed latency, and estimated dollars. Recent surveys of
test-time scaling likewise recommend reporting resource dimensions beyond final accuracy under
nominal budgets \citep{ttc_budget_survey,snell_ttc}. Our protocol follows that measurement view by
reporting componentwise realized resources and by auditing instrumentation artifacts that can alter
efficiency interpretation without altering correctness outcomes.

Recent systems-performance work reinforces the same scope for cloud and AI services: cloud
benchmarking methodology calls for explicit measurement and reporting practice; PEVA work on cloud
microservices studies uncertain latency--resource relationships under dynamic workloads; and PEVA
work on serverless ML treats public-cloud scaling and cost as part of the performance object
\citep{papadopoulos2021cloudmethodology,guo2023pobo,ali2025serverlessmltraining}. We adapt that
systems-performance view to closed-API LLM inference, where providers expose incomplete telemetry
and nominal budgets are not realized resource vectors.

\paragraph{Workload characterization and test-time controls}
Workload characterization asks whether compared systems face the same demand process and observable
resource contract. Adaptive compute-allocation work shows that how budget is spent can matter as
much as model size \citep{snell_ttc,adaptive_ttc_cpo,adactrl,ava}. Length-controlled and
budget-forcing approaches intervene directly on inference trajectories \citep{l1,s1,tale}. We treat
these mechanisms as workload controls in a shared nominal-$B$ environment and evaluate them under
common extraction and scoring conventions (Table~\ref{tab:baseline-fidelity}).

\paragraph{Discovery, consensus, and selection attribution}
Pass@$k$ analyses separate whether any sample is correct from whether a system selects it
\citep{passatk,ll_monkeys}. Majority vote is the standard same-model self-consistency (SC)
aggregator \citep{selfconsistency}; universal SC extends aggregation beyond extractable answers
\citep{universal_sc}. Adaptive and early-stopping SC variants exist
\citep{adaptive_consistency,esc,difficulty_adaptive_sc,reasc}. Best-of-$N$ and process verifiers
supply alternative selection signals \citep{self_certainty,letsverify}. Fixed-pool oracle-gap
diagnostics decompose recoverable mass and harm \citep{oracle_gap}. Cascades and routers choose
among models \citep{frugalgpt,hybrid_llm,routellm}. Relative to this literature, we keep the
generation--selection distinction as inherited, and instead ask how identical-pool plurality,
failure-mined deferral, and componentwise realized resources interact under one closed-API
multi-generator protocol. Pooled-4 targets a \emph{heterogeneous} four-generator pool---a different
axis from same-model SC---with one Azure$\times$GSM8K SC-$N{=}6$ budget cell as a supporting
control (Section~\ref{sec:sc-entitlement-control}).

\paragraph{Resource accounting in distributed/closed inference services}
Closed inference services expose partial telemetry: successful calls, retries, token fields,
latency, and billable cost may not align perfectly. This complicates cross-policy efficiency
comparisons even under matched nominal call budgets, and it creates a distributed-service
benchmarking problem rather than only an LLM-algorithm problem. Our instrumentation treatment aligns
with the performance-measurement principle that metrics must be interpreted at the layer where they
are observed, with reconstruction status disclosed explicitly.

\paragraph{Tree search and stopping}
Tree-of-Thoughts expands intermediate states under search heuristics \citep{tot}; classical
metareasoning frames stopping as value-of-computation \citep{metareasoning}. The evaluated Frontier
configuration does not test adaptive early stopping (Section~\ref{sec:method-frontier}). Historical
explorations outside the evaluated protocol are documented only in the supplementary material.
